% TODO checklist:
% - [x] Inspected src/workers/jobs_worker.py (_mark_cancelled, cancellation checks)
% - [x] Inspected db/redis_cache/job_store.py (cancel_job_atomic, check_cancel_flag)
% - [x] Inspected db/redis_cache/jobs_cache.py (cancel flag operations)
% - [x] Inspected src/routers/jobs.py (cancel_job endpoint)
% - [x] Analyzed Q&A.md Q39 for cancellation handling details

\section{Job Cancellation}
\label{sec:job-cancellation}

Job cancellation presents a coordination challenge in distributed systems: the cancel request originates from an API endpoint, but the job execution occurs in a separate worker process. The Cineca Agentic Platform implements a cooperative cancellation model using Redis flags for cross-process communication and Lua scripts for atomic state transitions.

\subsection{Cancellation Model}

The platform employs a \textit{cooperative cancellation} pattern where:

\begin{enumerate}
    \item A cancel flag is set in Redis by the API endpoint.
    \item The worker periodically checks this flag during job execution.
    \item If the flag is set, the worker gracefully terminates execution and transitions the job to CANCELLED status.
\end{enumerate}

This model ensures that jobs can complete their current unit of work before stopping, preventing resource leaks and data corruption from abrupt termination.

\subsection{Cancel Flag Mechanism}

The cancel flag is stored as a simple Redis key with the pattern \texttt{job:\{id\}:cancel}:

\begin{lstlisting}[language=Python, caption={Redis cancel flag operations in jobs\_cache.py.}]
def set_cancel_flag(job_id: str) -> bool:
    """Set cancellation flag for a job."""
    redis = get_redis_client()
    key = f"job:{job_id}:cancel"
    # Set flag with TTL matching job expiry
    return redis.setex(key, JOB_TTL_SECONDS, "1") == True

def check_cancel_flag(job_id: str) -> bool:
    """Check if job has been cancelled."""
    redis = get_redis_client()
    key = f"job:{job_id}:cancel"
    return redis.exists(key) == 1

def clear_cancel_flag(job_id: str) -> bool:
    """Clear cancellation flag (on job completion)."""
    redis = get_redis_client()
    key = f"job:{job_id}:cancel"
    return redis.delete(key) == 1
\end{lstlisting}

\subsection{Cancellation Flow}

Figure~\ref{fig:cancel-flow} illustrates the complete cancellation sequence across API and worker processes.

\begin{figure}[ht]
\centering
\begin{tikzpicture}[
    box/.style={rectangle, draw, minimum width=2cm, minimum height=0.6cm, font=\small},
    arrow/.style={->, >=stealth, thick},
    note/.style={font=\scriptsize, align=left}
]
    % Actors
    \node[box] (client) at (0,0) {Client};
    \node[box] (api) at (3,0) {API Server};
    \node[box] (redis) at (6,0) {Redis};
    \node[box] (worker) at (9,0) {Worker};
    
    % Timeline
    \draw[gray, dashed] (0,-0.5) -- (0,-5.5);
    \draw[gray, dashed] (3,-0.5) -- (3,-5.5);
    \draw[gray, dashed] (6,-0.5) -- (6,-5.5);
    \draw[gray, dashed] (9,-0.5) -- (9,-5.5);
    
    % Messages
    \draw[arrow] (0,-1) -- node[above, font=\scriptsize] {POST /jobs/\{id\}/cancel} (3,-1);
    \draw[arrow] (3,-1.5) -- node[above, font=\scriptsize] {SET cancel flag} (6,-1.5);
    \draw[arrow] (6,-2) -- node[above, font=\scriptsize] {OK} (3,-2);
    \draw[arrow] (3,-2.3) -- node[above, font=\scriptsize] {202 Accepted} (0,-2.3);
    
    % Worker polling
    \draw[arrow] (9,-3) -- node[above, font=\scriptsize] {GET cancel flag} (6,-3);
    \draw[arrow] (6,-3.3) -- node[above, font=\scriptsize] {flag = 1} (9,-3.3);
    
    % Status update
    \draw[arrow] (9,-4) -- node[above, font=\scriptsize] {status = cancelled} (6,-4);
    \draw[arrow] (9,-4.5) -- node[above, font=\scriptsize] {CLEAR cancel flag} (6,-4.5);
    
    % Notes
    \node[note, right=0.2cm of worker] at (9,-3.5) {Worker detects\\cancel request};
\end{tikzpicture}
\caption{Cancellation sequence diagram showing cross-process coordination.}
\label{fig:cancel-flow}
\end{figure}

\subsection{API Endpoint Implementation}

The cancel endpoint in \texttt{src/routers/jobs.py} handles cancel requests:

\begin{lstlisting}[language=Python, caption={Job cancellation API endpoint.}]
@router.post("/{job_id}/cancel")
async def cancel_job(
    job_id: str,
    user = Depends(get_current_principal),
):
    """Request cancellation of a running job."""
    _validate_job_id_uuid(job_id)
    
    # Get job and verify ownership
    job = await job_store.get(job_id)
    if not job:
        raise HTTPException(status_code=404, detail="Job not found")
    
    _require_owner_or_admin(user, job)
    
    # Check if already terminal
    if job.status.is_terminal:
        return {"job_id": job_id, "status": job.status.value, "cancelled": False}
    
    # Set cancel flag
    if settings.JOB_STORE_BACKEND == "redis":
        # Use atomic Lua script for Redis backend
        cancelled = await job_store.cancel_job_atomic(job_id)
    else:
        jobs_cache.set_cancel_flag(job_id)
        cancelled = True
    
    if cancelled:
        # Record cancellation event
        event_id = await event_store.get_next_event_id(job_id)
        event = SSEEvent(
            event_id=event_id,
            event_type="status",
            data={"to": "cancelled", "timestamp": datetime.utcnow().isoformat()}
        )
        await event_store.append(job_id, event, settings.SSE_RING_SIZE)
    
    return {
        "job_id": job_id,
        "cancel_requested": True,
        "status": "cancellation_pending"
    }
\end{lstlisting}

\subsection{Worker Cancellation Checking}

Workers check the cancel flag at multiple points during execution:

\begin{enumerate}
    \item \textbf{Before Execution}: Check immediately after dequeuing the job.
    \item \textbf{During Execution}: Check periodically within long-running operations.
    \item \textbf{After Execution}: Final check before marking as FINISHED.
\end{enumerate}

\begin{lstlisting}[language=Python, caption={Cancellation checking in worker execution flow.}]
async def _execute_job(self, job_id: str, db: Session):
    """Execute job with cancellation checking."""
    job_uuid = UUID(job_id)
    jobs_service = JobsService(db)
    
    # Check 1: Before starting
    if await asyncio.to_thread(jobs_cache.check_cancel_flag, job_id):
        logger.info(f"Job {job_id} was cancelled before execution")
        await self._mark_cancelled(job_uuid, jobs_service)
        return
    
    # ... transition to RUNNING ...
    
    # Execute with periodic cancellation checks
    result = await self._run_job_with_heartbeat(...)
    
    # Check 2: After execution
    if await asyncio.to_thread(jobs_cache.check_cancel_flag, job_id):
        logger.info(f"Job {job_id} was cancelled during execution")
        await self._mark_cancelled(job_uuid, jobs_service)
        return
    
    # ... transition to FINISHED ...
\end{lstlisting}

Within job execution, cancellation is checked at natural breakpoints:

\begin{lstlisting}[language=Python, caption={Periodic cancellation check within job execution.}]
async def _execute_long_running_job(self, job_id: str, payload: dict) -> dict:
    """Long-running job with cancellation checks per step."""
    steps = payload.get("steps", 10)
    
    for step in range(1, steps + 1):
        # Check before each step
        if await asyncio.to_thread(jobs_cache.check_cancel_flag, job_id):
            raise asyncio.CancelledError("Job cancelled")
        
        logger.info(f"Job {job_id}: step {step}/{steps}")
        await asyncio.sleep(3.0)  # Step duration
    
    return {"status": "completed", "steps_completed": steps}
\end{lstlisting}

\subsection{Atomic Cancellation with Lua Scripts}

For the Redis backend, atomic cancellation uses a Lua script to prevent race conditions between status checks and updates:

\begin{lstlisting}[language=Python, caption={Atomic cancellation using Lua script.}]
CANCEL_JOB_SCRIPT = """
-- KEYS[1] = job:{id}
-- ARGV[1] = timestamp
-- ARGV[2] = result JSON

local status = redis.call('HGET', KEYS[1], 'status')
if not status then
    return 'not_found'
end

-- Only cancel non-terminal jobs
if status == 'finished' or status == 'failed' or status == 'cancelled' then
    return 'already_terminal'
end

-- Atomic status update
redis.call('HSET', KEYS[1], 'status', 'cancelled')
redis.call('HSET', KEYS[1], 'updated_at', ARGV[1])
redis.call('HSET', KEYS[1], 'result', ARGV[2])

return 'cancelled'
"""

async def cancel_job_atomic(self, job_id: str) -> bool:
    """Atomically cancel using Lua CAS."""
    await self._ensure_scripts_loaded()
    redis = await get_async_redis()
    
    result = await redis.evalsha(
        self._cancel_job_sha,
        1,  # numkeys
        f"job:{job_id}",
        datetime.utcnow().isoformat(),
        json.dumps({"cancelled": True}),
    )
    
    result_str = result.decode("utf-8") if isinstance(result, bytes) else result
    
    if result_str == "cancelled":
        logger.info(f"Job {job_id} cancelled atomically")
        return True
    elif result_str == "already_terminal":
        logger.info(f"Job {job_id} already in terminal state")
        return False
    else:
        logger.warning(f"Job {job_id} not found for cancellation")
        return False
\end{lstlisting}

\subsection{Marking Jobs as Cancelled}

When cancellation is detected, the worker transitions the job to CANCELLED status:

\begin{lstlisting}[language=Python, caption={Worker method to mark job as cancelled.}]
async def _mark_cancelled(self, job_uuid: UUID, jobs_service: JobsService):
    """Mark job as cancelled in PostgreSQL."""
    try:
        # Try from queued state first
        job = jobs_service.transition_status(
            job_uuid, from_status="queued", to_status="cancelled"
        )

        if not job:
            # Try from running state
            job = jobs_service.transition_status(
                job_uuid, from_status="running", to_status="cancelled"
            )

        if job:
            # Log cancellation event for SSE
            jobs_service.append_event(
                job_uuid,
                event_type="status",
                event_data={
                    "to": "cancelled",
                    "timestamp": datetime.utcnow().isoformat()
                },
            )
            logger.info(f"Job {job_uuid} marked as cancelled")
            
    except Exception as e:
        logger.error(f"Failed to mark job {job_uuid} as cancelled: {e}")
\end{lstlisting}

\subsection{Cleanup Procedures}

After cancellation, cleanup ensures resources are properly released:

\begin{enumerate}
    \item \textbf{Cancel Flag Removal}: The cancel flag is cleared to prevent false positives on job restart scenarios.
    
    \item \textbf{Event Logging}: A terminal event is appended to the SSE stream.
    
    \item \textbf{Resource Release}: Any held resources (database connections, file handles) are released via context managers.
    
    \item \textbf{Metrics Recording}: Cancellation is recorded in Prometheus metrics.
\end{enumerate}

\begin{lstlisting}[language=Python, caption={Cancellation metrics tracking.}]
job_cancel_total = Counter(
    "job_cancel_total", "Total job cancellations",
    ["backend", "first_time"]  # first_time: true (202) or false (200)
)

job_cancel_duration_seconds = Histogram(
    "job_cancel_duration_seconds", "Job cancellation latency",
    ["backend"],
    buckets=(0.001, 0.005, 0.010, 0.025, 0.050, 0.100, 0.250, 0.500, 1.0)
)
\end{lstlisting}

\subsection{Cancellation Response Codes}

The API returns different HTTP status codes based on cancellation outcome:

\begin{table}[ht]
\centering
\caption{Cancellation API response codes.}
\label{tab:cancel-responses}
\small
\begin{tabular}{lll}
\hline
\textbf{Status Code} & \textbf{Condition} & \textbf{Response} \\
\hline
202 Accepted & First cancel request & \texttt{\{cancel\_requested: true\}} \\
200 OK & Already cancelled & \texttt{\{cancelled: true, status: "cancelled"\}} \\
200 OK & Already terminal & \texttt{\{cancelled: false, status: "finished"\}} \\
404 Not Found & Job doesn't exist & Error response \\
403 Forbidden & Not owner or admin & Error response \\
\hline
\end{tabular}
\end{table}

\subsection{Cancellation Timeout Considerations}

The cooperative cancellation model has inherent latency:

\begin{itemize}
    \item \textbf{Best Case}: Job checks cancel flag immediately, cancellation completes within milliseconds.
    
    \item \textbf{Typical Case}: Job is between check points, cancellation completes within one check interval (e.g., 0.5 seconds for demo jobs).
    
    \item \textbf{Worst Case}: Job is in an unchecked long operation, cancellation waits for operation completion.
\end{itemize}

For scenarios requiring immediate termination, operators can combine cooperative cancellation with process-level intervention (e.g., worker restart), though this may leave jobs in an inconsistent state.

